\section{Thảo luận}
\label{sec:discussion}

\subsection{Giá trị và trade-off của thiết kế}

Đóng góp thực tế của nguyên mẫu không phải là thay thế một nền tảng SIEM, mà là
liên kết một quyết định RBAC, luật phát hiện và bằng chứng trong cùng đầu ra.
Mỗi alert giữ định danh event, luật, evidence, điểm và severity; incident chỉ
gom các alert có cùng cặp người dùng--IP trong cửa sổ đã công bố. Vì vậy một
người phân tích có thể lần từ incident về event gốc mà không phải suy luận từ
một điểm số không giải thích được. Cách trình bày này phù hợp với mục đích quản
lý log và điều tra sau sự kiện, trong đó ngữ cảnh và khả năng truy vết cần được
duy trì \cite{Kent2006SP80092}.

Trade-off là tính linh hoạt thấp hơn. RBAC SQLite và các luật xác định giúp
kết quả ổn định, dễ test và tái lập, nhưng giả định role/permission trong seed
phản ánh đúng trạng thái ủy quyền tại thời điểm event. Tương tự, ngưỡng password
guessing, giờ làm việc, bonus và cửa sổ 300 giây là các quyết định cấu hình,
không phải tham số đã được hiệu chỉnh thống kê. Điều đó phù hợp với một pipeline
nhỏ gọn, song không loại bỏ nhu cầu quản trị policy liên tục trong một kiến trúc
zero-trust hay vận hành thực \cite{Rose2020SP800207}.

\subsection{Threats to validity}

\paragraph{Tính nội tại.} Dataset và \texttt{expected\_label} do nhóm tác giả
tạo để kiểm thử các kịch bản mục tiêu. Do đó precision, recall và F1 baseline
cho biết mức khớp của code với tập luật được thiết kế, không chứng minh khả năng
phát hiện trên dữ liệu độc lập. Một event có thể có nhiều finding ngữ cảnh, nên
không thể diễn giải số finding như số mẫu độc lập hoặc so sánh trực tiếp với
metric theo nhãn. Sai khác giữa timestamp, địa chỉ IP hoặc trạng thái RBAC trong
log và CSDL seed cũng có thể làm thay đổi kết quả mà không phản ánh một thay đổi
trong hành vi tấn công.

\paragraph{Tính xây dựng.} SQL injection được biểu diễn bằng một tập regex và
chỉ báo finding cho mẫu khớp đầu tiên; biến thể mã hóa, phân mảnh hoặc không có
trong mẫu có thể bị bỏ sót. Password guessing dựa trên số lần thất bại của đúng
cặp người dùng--IP; kẻ tấn công phân tán IP, tài khoản hay thời gian có thể
không đạt ngưỡng. Các tín hiệu IP mới, ngoài giờ, tài nguyên hiếm và từ chối
lặp lại là heuristic triage có thể tạo nhiễu cho công việc hợp lệ bất thường,
không phải kết luận về ý đồ hay thành công của tấn công. Chúng cần được xem như
đầu vào ưu tiên hóa để con người kiểm tra, thay vì một quyết định chặn truy cập.

\paragraph{Tính ngoại suy và vận hành.} Hai tập log đều tổng hợp, nhỏ và có
phong cách schema đồng nhất; chúng không phản ánh drift hành vi, NAT/proxy,
thiếu trường log, tải lớn, hay tác động riêng tư khi thu thập dữ liệu thực. Hệ
thống là hậu kiểm offline: không bảo vệ tính toàn vẹn/nguồn gốc log, không
tương quan đa nguồn, không thực hiện pháp chứng đầy đủ, phản ứng tự động hay
quyết định PDP thời gian thực. Các kiểm soát audit và authentication trong môi
trường triển khai cần được thiết kế độc lập, theo các yêu cầu tổ chức phù hợp
\cite{JTF2020SP80053,NIST2025SP80063B}.

\subsection{Hướng mở rộng có thể kiểm chứng}

Bước tiếp theo là đánh giá trên nhật ký đã ẩn danh, đại diện cho nhiều khoảng
thời gian và có nhãn/đánh giá của chuyên gia độc lập; khi đó cần báo cáo sai số,
độ trễ, chi phí triage và độ ổn định theo cấu hình, thay vì chỉ số đếm kịch bản.
Có thể bổ sung kiểm tra schema, provenance và correlation đa nguồn trước khi
đưa vào luồng sự kiện. Việc học baseline dài hạn hoặc mô hình học máy chỉ nên
được so sánh với heuristic hiện tại khi có dữ liệu đủ chất lượng, cách hiệu
chỉnh ngưỡng và cơ chế giải thích/kiểm toán rõ ràng. Những thay đổi này mở rộng
nguyên mẫu theo từng giả thuyết có thể đo được, thay vì suy diễn từ kết quả
trên dữ liệu tổng hợp hiện tại.
